\begin{frame}[t]{\ev{\tagO}250 attempts raised SWE-bench Lite coverage from 15.9\% to 56\%.}
\lead{DeepSeek-Coder-V2-Instruct; issues solved by at least one sample.}
{\fontsize{13}{16}\selectfont
\begin{tabularx}{\textwidth}{@{}Y r@{}}
\toprule
\textbf{Attempts per issue} & \textbf{Issue coverage}\\
\midrule
1 & 15.9\%\\
250 & \textbf{56\%}\\
\bottomrule
\end{tabularx}\par}
\vspace{.25cm}
\begin{ticks}
\item Coverage counts whether a correct candidate exists in the collection.
\item Production selection still needs reliable verification.
\item False acceptance costs can favor fewer than ten attempts.
\end{ticks}
\sources{Brown et al., arXiv:2407.21787; Stroebl et al., arXiv:2411.17501.}
\qadetail{Brown et al., Large Language Monkeys, arXiv:2407.21787: DeepSeek-Coder-V2-Instruct on SWE-bench Lite improves pass-at-least-once coverage from 15.9\% at one sample to 56\% at 250, a 40.1 percentage-point increase. This does not report production selection accuracy. Stroebl, Kapoor and Narayanan, The Limits of Inference Scaling Through Resampling, arXiv:2411.17501: when false-positive acceptance carries negative utility, the optimum is often fewer than ten attempts under their modeled conditions. This is not a universal branch-count recommendation. Neither paper measures our sandbox-fork effect.}
\note{[0:45] Brown and colleagues sampled DeepSeek-Coder-V2-Instruct on SWE-bench Lite. One sample solved about sixteen percent of issues. At 250 samples, at least one candidate solved fifty-six percent. That measures coverage of the collection. A production verifier must still recognize the useful candidate. Stroebl and colleagues show that false acceptance costs can make a small sampling budget preferable under their conditions. Our factory also had tests and reviews. Before connecting it to branching, we should separate what we observed, what we built, and what remains proposed.}
\end{frame}
